% !TeX root = ../main.tex
\section{问题描述与基本假设}

\subsection{系统与运营场景}
考虑由多个 O--D 对及沿线换电站组成的高速公路网络。车辆从入口进入后沿下游方向行驶，经必要的换电站补能并从出口离开。各站配备固定数量的电池和充电设备。一次换电取走满电电池并退回带有剩余电量的电池，退回电池经过充电后继续服务其他用户。因此，换电站的选择既影响车辆的后续行程，也影响各站电池 SOC 和充电需求。

系统服务预约用户和随机用户。预约用户提供 O--D、计划入口时间和入口 SOC，研究中的预约用户均具有满足续航条件的完整换电路径。尚未进入高速公路的用户由预约信息描述；进入高速公路后，使用其位置、SOC 和预计到站时间更新换电安排。已经到站且正在等待的请求保留原到站时刻、截止时刻和退回 SOC。随机用户直接到站提出请求，未来随机需求由当前预测信息描述。

\subsection{基本假设}
为描述上述运营过程，作如下假设。
\begin{enumerate}[label=\textbf{假设\arabic*.},leftmargin=4.4em]
  \item \textbf{道路与行驶。}车辆沿高速公路下游方向行驶，并按照给定换电站序完成行程。节点间距离和行驶能耗为已知参数，换电路径满足续航和出口最低 SOC 要求。研究场景中的入口 SOC 与初始路径相容，给定行驶能耗下已确定的换电安排在两次路径更新之间保持续航可行。
  \item \textbf{电池规格与库存。}可交换电池具有统一容量和兼容规格，服务就绪 SOC 归一化为 $1$。每次换电取走和退回各一块电池，站内电池数量保持不变。SOC $1$ 表示运营规定的服务目标。
  \item \textbf{充电过程。}充电效率为已知参数，电池 SOC 按输入电量线性变化。不同站点可以具有不同充电效率，各充电设备和站点分别具有给定功率上限。
  \item \textbf{换电操作。}换电视为瞬时操作。交付电池达到服务就绪 SOC，退回电池进入对应槽位并参加后续充电。槽位记录当时接入的电池状态。
  \item \textbf{信息与预测。}每轮可获得当前电池状态、等待请求及在途车辆的位置和 SOC。车辆到站时间采用当前 ETA 预测，未来随机需求和电价采用决策时刻可得的信息。同一用户沿下游方向的 ETA 按行驶顺序递增，预测信息随滚动过程更新。
  \item \textbf{到站顺序。}同一站点同类请求的原始到站时刻互不相同。同一时间段内可以存在多个请求，其先后次序由原始到站时刻确定。
\end{enumerate}

\subsection{运营规则}
预约请求优先于随机请求，同类请求按原到站时刻先后服务。每个请求自到站时刻起具有给定的最大等待时间，运营者可以在该窗口内选择服务时间，并为其匹配可用满电电池。预约在截止时刻仍未获得服务时，产生一次预约失败成本，该用户后续换电需求随之结束；随机请求在等待截止后流失。已完成的换电和行驶过程按实际记录保留，在线路径调整针对车辆后续的换电安排。换电安排实际发生变化时产生路径调整成本。

